 \section{Decoerenza Quantistica}
Fino ad ora abbiamo considerato il qubit come un sistema isolato che mantiene la sua informazione per un tempo indeterminato. Questo nella realtà non è propriamente vero, vi è un interazione con l'ambiente impossibile da evitare completamente, questo sarà poi tradotto in un entanglement con l'ambiente stesso, ma ne discuteremo successivamente. 

La questione della decoerenza rimane dcomunque aperta, questo può portare a dei problemi riguardanti lo stato del qubit nel tempo, ed al suo tempo in cui può essere processato (rispetto al passato dove il tempo era di diversi micro-secondi, adesso si può parlare anche di millisecondi o addirittura secondi. Questo dipende fortemente anche dalla natura del qubit, in questo stiamo considerando i Qubit Criogenici). Nel mondo classico, l'unico problema possibile è il bit flip, e quindi nel caso il dato corretto fosse 1 diventa 0 e vicevera. Mentre nel mondo quantum i principali problemi sono due:
\begin{itemize}
    \item \textbf{Bit-flip}: Che ugualmente al caso classico, uno stato in \(\ket{1}\) passa allo stato \(\ket{0}\) e viceversa. Questo corrisponde all'applicazione di una gate X di Pauli
    
    \item \textbf{Phase-flip (Dephasing)}: In questo caso le conseguenze andrebbero sulla base di Hadamard, dato che sarebbe come andare ad applicare una gate Z di Pauli
\end{itemize}

Per cercare di raggirare tali problematiche ci sono diversi modi e diverse tecniche che rientrano nell'ambito della \textbf{Quantum Error Correction} che si vedrà nei capitoli successivi.

